\scanpage{201}
\mspage{\unsure{57}}
$I'\mathbin{\unsure{-|\to}}I$ est ``transportable'' p. ex. si dans $I$ tout iso. est une identité, alors $u$ $W_0$-connexe implique que les fibres de $u$ sont $0$-connexes, comme on voit en \deleted{\illegible{1}} exprimant $d'')$ pour $\alpha$ flèches identiques. Les conditions envisagées sur $I$ sont satisfaites si $I$ est ordonnée. Mais $=$ si $I'$ et $I$ sont \inserted{\unsure{tous|toutes} deux} ordonnés, $u$ peut être $W_0$-connexe, sans être une $W_0$-équivalence universelle. Voici un exemple avec $I=\unsure{\mathbb{\Delta}|\Delta}^2$ :

\diagram{
\(\begin{aligned}
 I'&=\left\{\vcenter{\hbox{\begin{tikzcd}[ampersand replacement=\&,row sep=small,column sep=small]
 \&1'\arrow[r]\arrow[d]\&2\\
 0\arrow[r]\&1\&
 \end{tikzcd}}}\right\}&&\\[-0.7ex]
 &&&\downarrow\\[-0.7ex]
 I&=\left\{\vcenter{\hbox{\begin{tikzcd}[ampersand replacement=\&,column sep=small]
 0\arrow[r]\&1\arrow[r]\&2
 \end{tikzcd}}}\right\}&&
\end{aligned}\)
}

Les catégories bilocalisées de $I$ sont \inserted{les trois sommets} $\{0\},\{1\},\{2\}$, et les trois segments $\deleted{\unsure{\{}}[0,1\deleted{\unsure{\}}}],\deleted{\unsure{\{}}[1,2\deleted{\unsure{\}}}]$, et $[0,2]=I$, les images inverses sont respectivement

\diagram{
\(\bullet,\quad\)
\begin{tikzcd}[ampersand replacement=\&,row sep=small]
 \phantom{\cdot}\arrow[d]\\ \phantom{\cdot}
\end{tikzcd}
\(,\quad\bullet,\quad\)
\begin{tikzcd}[ampersand replacement=\&,row sep=small,column sep=small]
 \&\phantom{\cdot}\arrow[d]\\
 \phantom{\cdot}\arrow[r]\&\phantom{\cdot}
\end{tikzcd}
\(,\quad\)
\begin{tikzcd}[ampersand replacement=\&,row sep=small,column sep=small]
 \phantom{\cdot}\arrow[r]\arrow[d]\&\phantom{\cdot}\\
 \phantom{\cdot}\&
\end{tikzcd}
\(,\quad I'\deleted{\unsure{\Delta^2}},\)
}

elles sont toutes $0$-connexes, \unsure{ie} \deleted{\unsure{$=$}} $W_0$-asphériques, et même \deleted{\unsure{$W_\infty$-}} contractiles donc $W_\infty$-asphériques (donc $u$ est $=$ $W_\infty$-connexe, i.e. $W_\infty$-biasphérique). Mais ce n'est pas une $W_0$-équivalence universelle, car l'image inverse de $I'$ par $\mathbb{\Delta}^1\xrightarrow{\partial}\mathbb{\Delta}^2$, $\unsure{\partial|\delta}(0)=0$, $\partial(1)=2$,
% archive sheet number: 201
